\section{The settlement contract}
\label{sec:settlement}

The L1 component that holds the rollup state is \code{ShieldedPool}
\src{contracts/\allowbreak{}src/\allowbreak{}ShieldedPool.sol:63}. It stores a block counter and two
32-byte roots, delegates every validity question to an injected verifier, and
records the roots that a verified block statement names. It derives no root
itself, holds no nullifier set and no note data.
\src{contracts/\allowbreak{}src/\allowbreak{}ShieldedPool.sol:7-26,57-62}
This section states the state it holds, the transition \code{applyBlock}
enforces (in execution order), the event it emits, the fee path, the
deployment performed by \code{Deploy.s.sol}, and the calldata and transaction the
node sends to reach it. The statement vector $S$ and its decoding are those of
Def.~\ref{def:block-stmt} and Def.~\ref{def:block-decode}; the verifier's
binding of $S$ to the proof is specified in Section~7.

%---------------------------------------------------------------------------
\subsection{Contract state}
\label{subsec:pool-state}

\begin{definition}[Pool state]
\label{def:pool-state}
A \code{ShieldedPool} instance at address $A$ on a chain with identifier
$\chi$ is the tuple
\[
  \mathcal{P} \;=\; \big(\, V,\ f \;;\ N,\ R,\ R^{\mathsf{nf}} \,\big),
\]
with the following components.
\begin{center}
\small
\begin{tabularx}{\linewidth}{@{}l l l l X@{}}
\toprule
Symbol & Solidity name & Type & Kind & Written by \\
\midrule
$V$ & \code{verifier} & \code{IWhirVerifier} & \code{immutable} & constructor only \src{contracts/\allowbreak{}src/\allowbreak{}ShieldedPool.sol:68,131} \\
$f$ & \code{feeRecipient} & \code{address} & storage & constructor only \src{contracts/\allowbreak{}src/\allowbreak{}ShieldedPool.sol:89,132} \\
$N$ & \code{blockNumber} & \code{uint256} & storage & \code{applyBlock} (initially $0$) \src{contracts/\allowbreak{}src/\allowbreak{}ShieldedPool.sol:71,168-170} \\
$R$ & \code{currentRoot} & \code{bytes32} & storage & constructor, \code{applyBlock} \src{contracts/\allowbreak{}src/\allowbreak{}ShieldedPool.sol:74,133,165} \\
$R^{\mathsf{nf}}$ & \code{currentNullifierRoot} & \code{bytes32} & storage & constructor, \code{applyBlock} \src{contracts/\allowbreak{}src/\allowbreak{}ShieldedPool.sol:77,134-135,166} \\
$E^{\mathsf{nf}}$ & \code{EMPTY\_NULLIFIER\_ROOT} & \code{bytes32} & \code{constant} & --- \src{contracts/\allowbreak{}src/\allowbreak{}ShieldedPool.sol:85-86} \\
\bottomrule
\end{tabularx}
\end{center}
All five state variables and the constant are \code{public}, so each has an
auto-generated getter. The view selectors for $N$, $R$, $R^{\mathsf{nf}}$
are pinned off-chain as \code{0x57e871e7}, \code{0xfdab463d},
\code{0x222d1bed}. \src{contracts/\allowbreak{}test/\allowbreak{}SelectorPins.t.sol:14-18}
\end{definition}

\begin{property}[Write surface]
\label{prop:pool-writes}
Besides the constructor, the contract declares exactly two external
functions, \code{applyBlock} and \code{withdrawFees}
\src{contracts/\allowbreak{}src/\allowbreak{}ShieldedPool.sol:144,182}. Only \code{applyBlock} writes
$(N,R,R^{\mathsf{nf}})$; nothing writes $f$ after construction, and $V$ is
immutable. There is no owner, no pause, no upgrade path and no proxy. No
function is \code{payable}, and there is no \code{receive} or \code{fallback}
function. \src{contracts/\allowbreak{}src/\allowbreak{}ShieldedPool.sol:63-190}
\end{property}

%---------------------------------------------------------------------------
\subsection{Construction}
\label{subsec:pool-ctor}

\begin{definition}[Constructor]
\label{def:pool-ctor}
\code{constructor}$(V, f, r_0, r^{\mathsf{nf}}_0)$ enforces
$V \ne 0$ and $f \ne 0$ (string reverts), then sets
\begin{align}
  R &:= r_0, \\
  R^{\mathsf{nf}} &:= \begin{cases} E^{\mathsf{nf}} & \text{if } r^{\mathsf{nf}}_0 = 0^{32},\\
                                  r^{\mathsf{nf}}_0 & \text{otherwise,}\end{cases}
\end{align}
and leaves $N = 0$.
\src{contracts/\allowbreak{}src/\allowbreak{}ShieldedPool.sol:123-136}
No check is applied to $r_0$ or to a nonzero $r^{\mathsf{nf}}_0$, and $V$ is not
checked to have code. A wrong genesis pair is not rejected at construction; it
makes every block whose statement was witnessed against the true genesis fail
the continuity checks of Def.~\ref{def:apply-block}.
\src{contracts/\allowbreak{}src/\allowbreak{}ShieldedPool.sol:120-122}
\end{definition}

\begin{specdelta}[Empty nullifier root]
\label{delta:pool-empty-nf}
The constant is
\begin{Verbatim}[fontsize=\small]
EMPTY_NULLIFIER_ROOT =
  0x83265b7a2c40cc6eca237b6016b764606335cd0c6ff0732470e20a484d5a7329
\end{Verbatim}
documented as the Poseidon2 empty-subtree chain at depth
$\code{NULLIFIER\_TREE\_DEPTH}=96$
\src{contracts/\allowbreak{}src/\allowbreak{}ShieldedPool.sol:79-86}
\src{crates/\allowbreak{}shielded/\allowbreak{}src/\allowbreak{}nullifier\_tree.rs:64}, and it equals the
\code{nullifier\_before\_hex} of the settled-block vector
\src{contracts/\allowbreak{}test/\allowbreak{}vectors/\allowbreak{}block\_vectors.json:4}
(see Section~\ref{subsec:ss-constants}). The constructor documentation instead
states that a zero argument makes the contract \emph{compute} a depth-256
Keccak empty root at deploy time
\src{contracts/\allowbreak{}src/\allowbreak{}ShieldedPool.sol:114-118}; the code performs no hashing and
selects the pinned constant \src{contracts/\allowbreak{}src/\allowbreak{}ShieldedPool.sol:134-135}
\finding{I-02}.
\end{specdelta}

%---------------------------------------------------------------------------
\subsection{\code{applyBlock}: the state transition}
\label{subsec:pool-apply}

\begin{definition}[\code{applyBlock}]
\label{def:apply-block}
\code{applyBlock(uint256[] calldata S, bytes calldata }$\pi$\code{)} is
\code{external}, non-\code{payable}, and has no caller restriction.
\src{contracts/\allowbreak{}src/\allowbreak{}ShieldedPool.sol:144}
It executes the following steps in order; any failing step reverts the whole
call, so no step's effect persists unless all steps succeed.
\begin{center}
\small
\begin{tabularx}{\linewidth}{@{}r X l@{}}
\toprule
\# & Step & Source \\
\midrule
1 & $b := V.\code{verify}(S,\pi)$. The interface declares \code{verify}
    \code{view}, so the call is a static call; a revert inside $V$ propagates.
    If $b=\code{false}$, revert \code{NotVerified()}.
  & \src{ShieldedPool.sol:145} \src{IWhirVerifier.sol:21-23} \\
2 & $B := \code{decode}(S)$ after copying $S$ from calldata to memory; all
    checks of Def.~\ref{def:block-decode} apply (string reverts).
  & \src{ShieldedPool.sol:147-149} \\
3 & require $B.\code{rootBefore} = R$, else revert
    \code{RootMismatch}$(R,\ B.\code{rootBefore})$.
  & \src{ShieldedPool.sol:152-154} \\
4 & require $B.\code{nullifierBefore} = R^{\mathsf{nf}}$, else revert
    \code{NullifierRootMismatch}$(R^{\mathsf{nf}},\ B.\code{nullifierBefore})$.
  & \src{ShieldedPool.sol:155-157} \\
5 & $R := B.\code{rootAfter}$;\ $R^{\mathsf{nf}} := B.\code{nullifierAfter}$.
  & \src{ShieldedPool.sol:165-166} \\
6 & $N := N+1$ (in an \code{unchecked} block).
  & \src{ShieldedPool.sol:168-170} \\
7 & emit \code{BlockApplied}$(N, B.\code{rootBefore}, B.\code{rootAfter},
    B.\code{nullifierBefore}, B.\code{nullifierAfter}, B.\code{totalFee})$
    with the incremented $N$.
  & \src{ShieldedPool.sol:171-178} \\
\bottomrule
\end{tabularx}
\end{center}
\end{definition}

The guard portion of the function, verbatim:
\begin{Verbatim}[fontsize=\small]
if (!verifier.verify(statement, proof)) revert NotVerified();
BlockStatement.Block memory block_ = statement.decode();
if (block_.rootBefore != currentRoot) {
    revert RootMismatch(currentRoot, block_.rootBefore);
}
if (block_.nullifierBefore != currentNullifierRoot) {
    revert NullifierRootMismatch(currentNullifierRoot, block_.nullifierBefore);
}
\end{Verbatim}
\src{contracts/\allowbreak{}src/\allowbreak{}ShieldedPool.sol:145-157}

Writing $D_j(S) = \code{digestFromLimbs}(S,\ 1+2n+16j)$ for $j=0,\dots,4$
(Def.~\ref{def:digest-from-limbs}), the enforced relation is:

\begin{definition}[Settlement relation]
\label{def:settle-rel}
For a pool state $(V,f;N,R,R^{\mathsf{nf}})$, \code{applyBlock}$(S,\pi)$
succeeds if and only if
\begin{align}
  &V.\code{verify}(S,\pi) = \code{true} \text{ (without reverting)}, \label{eq:settle-v}\\
  &\code{decode}(S) \text{ succeeds, i.e. } 1\le S[0]=n\le 2^{16}-1,\ |S| = 81+6n,
     \text{ all decoded limbs } \le \code{0xffff}, \label{eq:settle-d}\\
  &D_1(S) = R \quad\text{and}\quad D_3(S) = R^{\mathsf{nf}}, \label{eq:settle-c}
\end{align}
and on success the post-state is
\begin{equation}
  (N,\ R,\ R^{\mathsf{nf}}) \;\mapsto\; (N+1,\ D_2(S),\ D_4(S)).
\end{equation}
\src{contracts/\allowbreak{}src/\allowbreak{}ShieldedPool.sol:144-179}
\src{contracts/\allowbreak{}src/\allowbreak{}BlockStatement.sol:79-112}
No other input, state variable or environment value (\code{msg.sender},
\code{tx.origin}, \code{block.chainid}, \code{address(this)},
\code{block.number}, \code{msg.value}) enters the relation.
\end{definition}

\begin{property}[Decoded but unused fields]
\label{prop:apply-unused}
Of the decoded \code{Block}, \code{numTransfers} and
\code{statementRoot} $=D_0(S)$ are not read by \code{applyBlock}, and
\code{totalFee} is used only as an event argument.
\src{contracts/\allowbreak{}src/\allowbreak{}ShieldedPool.sol:149-178}
The header limbs $S[1],\dots,S[2n]$ are never read on L1
(Def.~\ref{def:block-decode}). The fold root is therefore constrained on L1
only through \eqref{eq:settle-v} \finding{L-04}.
\end{property}

\begin{property}[Caller]
\label{prop:apply-caller}
\code{applyBlock} contains no \code{msg.sender} comparison, modifier or
allow-list; every account may call it, and the only gate is
\eqref{eq:settle-v}--\eqref{eq:settle-c}.
\src{contracts/\allowbreak{}src/\allowbreak{}ShieldedPool.sol:144-179}
The adequacy of this gate is examined in \finding{H-01}.
\end{property}

\begin{property}[Statement context]
\label{prop:apply-context}
The statement layout (Def.~\ref{def:block-stmt}) carries the transfer count,
the shape header, the fold root, the four chain endpoints and the fees, and
no chain identifier, pool address, block height or other per-deployment value.
\src{contracts/\allowbreak{}src/\allowbreak{}BlockStatement.sol:11-15,79-112}
The only binding between a pair $(S,\pi)$ and a particular pool instance is
\eqref{eq:settle-c}, i.e.\ equality of $D_1(S), D_3(S)$ with that instance's
current roots \finding{M-01}.
\end{property}

\begin{property}[Re-submission on the same instance]
\label{prop:apply-resubmit}
After $(S,\pi)$ has been applied, $R = D_2(S)$ and $R^{\mathsf{nf}} = D_4(S)$.
A second \code{applyBlock}$(S,\pi)$ on the same instance passes
\eqref{eq:settle-c} only if $D_2(S)=D_1(S)$ and $D_4(S)=D_3(S)$, i.e.\ only for a
block that leaves both roots unchanged.
\src{contracts/\allowbreak{}src/\allowbreak{}ShieldedPool.sol:152-157,165-166}
\end{property}

\begin{property}[Check order and cost]
\label{prop:apply-order}
Steps 2--4, which read only $S$ and two storage slots, run after step 1. A
statement that is malformed (fails \eqref{eq:settle-d}) or stale (fails
\eqref{eq:settle-c}) is therefore rejected only after the verifier has
executed. The ordering is documented as deliberate.
\src{contracts/\allowbreak{}src/\allowbreak{}ShieldedPool.sol:137-145}
The node submits each settlement with a gas ceiling of $2\cdot 10^{10}$,
described as matching the dev-chain configuration because the verifier's cost
is data-dependent.
\src{crates/\allowbreak{}node/\allowbreak{}src/\allowbreak{}settlement.rs:28-31}
\end{property}

\begin{property}[Fee handling in the transition]
\label{prop:apply-fee}
$\code{totalFee} = \sum_{c=1}^{n}\code{u64FromLimbs}(S,\ 81+2n+4(c-1))$ is
accumulated in checked \code{uint256} arithmetic
\src{contracts/\allowbreak{}src/\allowbreak{}BlockStatement.sol:104-107}; \code{applyBlock} moves no
value and credits no account with it, it only emits it.
\src{contracts/\allowbreak{}src/\allowbreak{}ShieldedPool.sol:171-178}
The contract does not check fee correctness; it relies on the transfer
relation's value conservation.
\src{contracts/\allowbreak{}src/\allowbreak{}ShieldedPool.sol:59-62}
\end{property}

%---------------------------------------------------------------------------
\subsection{The \code{BlockApplied} event}
\label{subsec:pool-event}

\begin{definition}[Event]
\label{def:pool-event}
\begin{Verbatim}[fontsize=\small]
event BlockApplied(
    uint256 indexed blockNumber,
    bytes32 rootBefore, bytes32 rootAfter,
    bytes32 nullifierRootBefore, bytes32 nullifierRootAfter,
    uint256 fee);
\end{Verbatim}
\src{contracts/\allowbreak{}src/\allowbreak{}ShieldedPool.sol:91-99}
The log has topic$_0$ = the event signature hash, topic$_1 = N$ (the
post-increment value, so the first applied block logs $1$), and five data
words $D_1(S), D_2(S), D_3(S), D_4(S), \code{totalFee}$.
\src{contracts/\allowbreak{}src/\allowbreak{}ShieldedPool.sol:168-178}
\end{definition}

\begin{property}[Event content]
\label{prop:pool-event-content}
The event omits the fold root $D_0(S)$, the transfer count $n$, the shape
header, the per-transfer fees, and the submitting account. Neither the event
nor $S$ carries per-transfer nullifiers or output commitments: on L1 these
appear only as absorbed into the fold root (Def.~\ref{def:block-fold})
\finding{L-04}.
\src{contracts/\allowbreak{}src/\allowbreak{}ShieldedPool.sol:91-99,171-178}
\src{contracts/\allowbreak{}src/\allowbreak{}BlockStatement.sol:17-24}
\end{property}

%---------------------------------------------------------------------------
\subsection{\code{withdrawFees}}
\label{subsec:pool-fees}

\begin{definition}[\code{withdrawFees}]
\label{def:withdraw-fees}
\code{withdrawFees()} requires $\code{msg.sender} = f$, else reverts
\code{NotFeeRecipient()}; with $a := \code{address(this).balance}$, if $a>0$ it
performs \code{f.call\{value: a\}("")} and requires success.
\src{contracts/\allowbreak{}src/\allowbreak{}ShieldedPool.sol:181-189}
\end{definition}

\begin{specdelta}[Fee accrual]
\label{delta:pool-fees}
$f$ is documented as the recipient of ``collected fees''
\src{contracts/\allowbreak{}src/\allowbreak{}ShieldedPool.sol:88-89}. By Property~\ref{prop:pool-writes}
and Property~\ref{prop:apply-fee}, no function of the contract accepts value
and \code{applyBlock} credits nothing, so no code path of the system funds the
balance that \code{withdrawFees} pays out; the shielded fee $\varphi$ of each
transfer has no on-chain counterpart in this contract \finding{I-04}.
\end{specdelta}

%---------------------------------------------------------------------------
\subsection{Deployment}
\label{subsec:pool-deploy}

\begin{definition}[Deployment script]
\label{def:deploy}
\code{Deploy.run()} performs, in order:
\begin{center}
\small
\begin{tabularx}{\linewidth}{@{}r X l@{}}
\toprule
\# & Action & Source \\
\midrule
1 & Read the genesis file named by \code{GENESIS\_FILE}, defaulting to
    \code{test/vectors/block\_genesis.json}; $r_0 :=$ its
    \code{.genesis\_root\_hex} as \code{bytes32}.
  & \src{Deploy.s.sol:28-32} \\
2 & $\mathit{dep} := \code{msg.sender}$ of the script (the broadcast account).
  & \src{Deploy.s.sol:34-36} \\
3 & Broadcast \code{new TerminalWeight()} $=T$.
  & \src{Deploy.s.sol:37-40} \\
4 & Broadcast \code{new WhirVerifier}$(T)$ $=W$.
  & \src{Deploy.s.sol:41} \\
5 & Broadcast \code{new ShieldedPool}$(W,\ \mathit{dep},\ r_0,\ 0^{32})$, so
    $f=\mathit{dep}$ and $R^{\mathsf{nf}} = E^{\mathsf{nf}}$.
  & \src{Deploy.s.sol:42-46} \\
6 & Write \code{deployments/local.json} with keys \code{verifier},
    \code{terminalWeight}, \code{pool}, \code{deployer}, \code{chainId}.
  & \src{Deploy.s.sol:48-57} \\
\bottomrule
\end{tabularx}
\end{center}
The documented invocation targets an \code{anvil} chain with chain id
$31337$ and the dev chain's default funded key.
\src{contracts/\allowbreak{}script/\allowbreak{}Deploy.s.sol:4-6}
\end{definition}

\begin{property}[Verifier instance]
\label{prop:deploy-verifier}
The constructor of \code{WhirVerifier} rejects a zero or code-less
\code{terminalWeight} and pins its address and code hash, which are re-checked
before every satellite call.
\src{contracts/\allowbreak{}src/\allowbreak{}verifier/\allowbreak{}WhirVerifier.sol:116-132,895}
It takes no configuration argument; \code{verify} decodes the CONFIG section
from the caller-supplied bundle $\pi$.
\src{contracts/\allowbreak{}src/\allowbreak{}verifier/\allowbreak{}WhirVerifier.sol:209-213,262-266}
The script does not deploy \code{WhirVerifierV6}, whose constructor pins a
CONFIG digest and length
\src{contracts/\allowbreak{}src/\allowbreak{}verifier/\allowbreak{}WhirVerifierV6.sol:41-47,61-75}; consequently
the pool's immutable $V$ is the bare engine. See Section~7 and
\finding{V-01}.
\end{property}

\begin{specdelta}[Explicit genesis nullifier root]
The script comments that ``the genesis file may pin one explicitly''; the code
passes the literal $0^{32}$ and reads no nullifier root from the genesis
file. \src{contracts/\allowbreak{}script/\allowbreak{}Deploy.s.sol:42-45}
\end{specdelta}

%---------------------------------------------------------------------------
\subsection{Settlement calldata}
\label{subsec:pool-calldata}

The node constructs the call by hand.
\src{crates/\allowbreak{}node/\allowbreak{}src/\allowbreak{}settlement.rs:1-8}

\begin{definition}[Word encodings]
\label{def:settle-enc}
For $x \in [0,2^{64})$, $\mathsf{w}(x)\in\{0,1\}^{256}$ is the 32-byte big-endian
word with the top 24 bytes zero
\src{crates/\allowbreak{}node/\allowbreak{}src/\allowbreak{}settlement.rs:93-98}. For a sequence
$s = (s_0,\dots,s_{L-1})$ of \code{u64} and a byte string $b$:
\begin{align}
  \mathsf{arr}(s) &= \mathsf{w}(L)\,\|\,\mathsf{w}(s_0)\,\|\cdots\|\,\mathsf{w}(s_{L-1}),
  \\
  \mathsf{byt}(b) &= \mathsf{w}(|b|)\,\|\,b\,\|\,0^{8\rho},\qquad
  \rho = (32 - |b| \bmod 32) \bmod 32 .
\end{align}
\src{crates/\allowbreak{}node/\allowbreak{}src/\allowbreak{}settlement.rs:100-116}
\end{definition}

\begin{definition}[\code{applyBlock} calldata]
\label{def:settle-calldata}
\begin{equation}
  \mathsf{cd}(S,\pi) \;=\; \code{0x0cb000b5}\ \|\ \mathsf{w}(\code{0x40})\ \|\
  \mathsf{w}\big(\code{0x40} + 32(L+1)\big)\ \|\ \mathsf{arr}(S)\ \|\ \mathsf{byt}(\pi),
  \qquad L = |S|,
\end{equation}
of length $4 + 64 + 32(L+1) + 32 + |\pi| + \rho$.
\src{crates/\allowbreak{}node/\allowbreak{}src/\allowbreak{}settlement.rs:118-133}
The two head words are the ABI offsets of the two dynamic tails measured from
the end of the selector; this is the standard ABI encoding of
\code{(uint256[],bytes)}. The selector is
$\code{bytes4}(\code{keccak256}(\code{"applyBlock(uint256[],bytes)"}))$, pinned
by a test that compares it with the EVM's own \code{keccak256}.
\src{crates/\allowbreak{}node/\allowbreak{}src/\allowbreak{}settlement.rs:24-26}
\src{contracts/\allowbreak{}test/\allowbreak{}SelectorPins.t.sol:10-12}
For a block statement, $L = 81+6n$ (Def.~\ref{def:block-stmt}).
\end{definition}

\begin{property}[Statement words]
The node takes $S$ as the canonical representatives of the block statement's
field elements, widened to \code{u64}, and $\pi$ as the settlement bundle
encoded from the same recursion artifact and statement.
\src{crates/\allowbreak{}node/\allowbreak{}src/\allowbreak{}main.rs:444-456}
Every word of $\mathsf{arr}(S)$ is thus $< p < 2^{32}$; the encoder cannot
express \code{uint256} values $\ge 2^{64}$, and none are needed.
\end{property}

%---------------------------------------------------------------------------
\subsection{Submission}
\label{subsec:pool-submit}

\begin{definition}[Settlement transaction]
\label{def:settle-tx}
\code{SettlementSender::settle}$(S,\pi)$ issues one JSON-RPC
\code{eth\_sendTransaction} with the object
\begin{center}
\small
\begin{tabular}{@{}l l l@{}}
\toprule
Field & Value & Source \\
\midrule
\code{from} & manifest \code{deployer} & \src{settlement.rs:201} \\
\code{to} & manifest \code{pool} & \src{settlement.rs:202} \\
\code{data} & \code{0x}\,$\|$\,hex$(\mathsf{cd}(S,\pi))$ & \src{settlement.rs:199,203} \\
\code{gas} & \code{0x4a817c800} $= 2\cdot 10^{10}$ & \src{settlement.rs:31,204} \\
\bottomrule
\end{tabular}
\end{center}
and no \code{value}, \code{nonce}, fee or \code{chainId} field; signing and the
remaining fields are left to the RPC endpoint, and the node holds no key.
\src{crates/\allowbreak{}node/\allowbreak{}src/\allowbreak{}settlement.rs:10-17,198-216}
On a dev chain the \code{from} account is an unlocked account of the node
(\code{anvil}).
\src{crates/\allowbreak{}node/\allowbreak{}src/\allowbreak{}settlement.rs:12-17}
\end{definition}

\begin{definition}[Receipt wait]
\label{def:settle-wait}
After the send, the sender polls \code{eth\_getTransactionReceipt} with a
$2$\,s sleep between polls. A non-null receipt with
\code{status} $=$ \code{"0x1"} returns the hash; any other status (absent
status read as \code{"0x0"}) returns \code{Reverted}. The $600$\,s deadline is
evaluated only after a poll that returned a null receipt.
\src{crates/\allowbreak{}node/\allowbreak{}src/\allowbreak{}settlement.rs:148-156,218-250}
The HTTP client is \code{reqwest::Client::new()} with no request timeout
configured, and \code{settle} is awaited inline by the node's single command
actor, which processes commands sequentially.
\src{crates/\allowbreak{}node/\allowbreak{}src/\allowbreak{}settlement.rs:150}
\src{crates/\allowbreak{}node/\allowbreak{}src/\allowbreak{}main.rs:464-476,507-509,780-781}
See \finding{I-05}.
\end{definition}

\begin{property}[Manifest and retention]
The node reads the manifest from \code{NODE\_MANIFEST} (default
\code{contracts/deployments/local.json}) and the RPC URL from
\code{NODE\_RPC\_URL} (default \code{http://127.0.0.1:8545})
\src{crates/\allowbreak{}node/\allowbreak{}src/\allowbreak{}main.rs:125-126}; the \code{Manifest} type
deserializes only \code{pool} and \code{deployer}, so the \code{chainId} the
script writes is not consulted
\src{crates/\allowbreak{}node/\allowbreak{}src/\allowbreak{}settlement.rs:70-77}. A pending
$(S,\pi)$ is cleared only on success and is retained on any error.
\src{crates/\allowbreak{}node/\allowbreak{}src/\allowbreak{}main.rs:464-476}
\end{property}

\begin{specdelta}[Settlement-side documentation]
\label{delta:settle-docs}
(i) \code{Manifest.pool} is documented as ``the \code{ShieldedPool} proxy
address'' and the manifest as written by \code{contracts/scripts/deploy\_local.mjs}
\src{crates/\allowbreak{}node/\allowbreak{}src/\allowbreak{}settlement.rs:72,80-81}; the deployment
deploys no proxy and is \code{Deploy.s.sol}, and no \code{deploy\_local.mjs}
exists under \code{contracts/scripts}.
\src{contracts/\allowbreak{}script/\allowbreak{}Deploy.s.sol:13,45}
(ii) The \code{settle} documentation states that re-sending a block whose first
attempt landed ``would double-apply nullifiers''
\src{crates/\allowbreak{}node/\allowbreak{}src/\allowbreak{}settlement.rs:190-193}; by
Property~\ref{prop:apply-resubmit}, the contract rejects such a re-send at
step 3 or 4 unless the block leaves both roots unchanged.
\end{specdelta}
